% 原书分章 PDF 第 2—6 页；章首导言与第 20.1—20.3 节。第 1 页第六编扉页由总审处理。
\chapter{纵向网络相依数据中的因果推断}
\label{ch:longitudinal-network-dependent-data}
\chapterauthors{Oleg Sofrygin 与 Mark J. van der Laan}
\chapteraffiliations{%
O. Sofrygin（通讯作者）：美国加利福尼亚大学伯克利分校生物统计学部，101 Haviland Hall, \#7358, Berkeley, CA 94720, USA；电子邮箱：\href{mailto:sofrygin@berkeley.edu}{sofrygin@berkeley.edu}。\\
M. J. van der Laan：美国加利福尼亚大学伯克利分校生物统计学部及统计学系，101 Haviland Hall, \#7358, Berkeley, CA 94720, USA；电子邮箱：\href{mailto:laan@berkeley.edu}{laan@berkeley.edu}。\\
本章 DOI：\url{https://doi.org/10.1007/978-3-319-65304-4_20}。}

现有因果推断文献大多关注这样的问题：观察 $n$ 个在因果上互不相连的独立单位的纵向数据，据此研究单时点或多时点干预对结局的效应。此时，因果效应被定义为分配给某个单位的干预对该单位结局的效应之期望，并假定施加于其他单位的干预不会对该单位的结局产生因果效应。因此，因果模型只需考虑单位自身数据结构各分量之间的因果关系。统计推断则建立在以下假设上：这 $n$ 个数据结构是某个随机变量的 $n$ 个独立实现值。然而，在许多整群随机试验（CRTs），或者只涉及少数社区的观察性研究中，独立单位的数量不足以支持基于极限分布的统计推断。

\begin{sourcebox}
在一种极端但并不少见的情形中，我们可能只观察到一个由因果相连的个体组成的社区。这时，是否仍能统计评估分配于社区层面的干预，对社区层面结局，例如个体结局的平均值，所产生的因果效应？这正是本章要回答的问题。显然，为定义所需的因果量，必须采用纳入所有单位的因果模型；而且，应在不依赖独立单位数量的渐近理论的前提下，建立这些因果量在最小假设下的可识别性。
\end{sourcebox}

如前所述，当数据采集自一个由相互连接的单位组成的总体时，经常采用的单位间独立假设通常并不成立，因为网络交互往往使某个单位的暴露影响其他相连单位的结局。考虑单时点暴露的情形：我们收集 $N$ 个相连单位的基线协变量、暴露和结局数据。可以预期，任意两个相连单位之间的交互，都可能使一个单位的暴露影响另一个单位的结局；这种现象通常称为个体间干扰（interference）或溢出效应（spillover）（Sobel，2006；Hudgens 和 Halloran，2008）。

虽然过去许多研究为了简化问题，错误地假定不存在个体间干扰，\footnote{关于错误忽略个体间干扰所带来的后果，参见第 21.7 节所列文献。}但在过去十年中，研究网络相依数据（network-dependent data）中因果效应估计的文献也日益增多。许多研究试图回答社交网络在公共卫生各方面所起作用的问题。例如，Christakis 等在一系列研究中，使用由社交网络相连的研究对象的观察性数据，估计肥胖、吸烟及其他多种结局的传染因果效应，发现许多状况都存在社会传染（social contagion）（Christakis 和 Fowler，2007，2013）。其中一项研究发现，即使控制了所有已测量混杂因素，个体每多一位肥胖朋友，其变得肥胖的风险也会增加（Christakis 和 Fowler，2007）。不过，这些研究使用的统计方法受到了审视：它们可能没有考虑观测单位间的网络相依性，因而低估标准误，使推断不够保守（Lyons，2010）；效应估计也可能存在偏差，其原因包括模型误设（Lyons，2010；VanderWeele 和 An，2013）、网络诱导的同质偏好（homophily）（Shalizi 和 Thomas，2011），以及环境因素造成的未测量混杂（Shalizi 和 Thomas，2011）。

\section{建模方法}

我们的建模方法有一个重要组成部分，就是纳入网络信息：对一个由 $N$ 个单位组成的有限总体中的每个单位 $i$，用集合 $F_i\subset\{1,\ldots,N\}$ 描述可能向该单位提供输入的其他单位。这样，便可为这组单位提出一个结构方程模型：除外生误差外，单位 $i$ 在时点 $t$ 的观测数据节点，仅受到 $F_i$ 中单位的观测数据的因果影响。集合 $F_i$ 只包含在各时点 $t$ 直接影响单位 $i$ 数据的直接朋友。这一结构方程模型可用一个包含全部 $N$ 个单位的所谓因果图（causal graph）可视化；我们可以将其称为网络。

我们对结构方程模型中外生误差的假设，对应于以下假设：给定所有单位的过去，时点 $t$ 各单位的数据节点具有序贯条件独立性。也就是说，给定所有单位最近的过去，包括网络信息，下一个时点各单位的数据彼此独立。单位 $i$ 的朋友集合 $F_i$ 越小，因果图中每个节点的入边就越少，从而能为所需量的靶向估计提供更大的有效样本量（effective sample size）。虽然这些因果图允许单位之间以复杂方式相依，但如果 $F_i$ 的大小关于 $N$ 一致有界，再结合外生误差独立的假设，全部 $N$ 个单位的数据似然便允许我们以单位数 $N$，而非社区数，例如 1，来驱动统计推断。此外，我们最近的工作将这些正式的渐近结果推广到允许 $F_i$ 的大小随 $N$ 增长的情形（Ogburn 等，2017）。

我们将运用因果效应靶向学习的路线图，定义并解决这一估计问题。首先，定义结构因果模型（structural causal model，SCM；Pearl，2009a），确定每个节点如何由其父数据节点及外生变量的函数生成。然后，通过对各单位处理节点施加随机干预（stochastic intervention），定义感兴趣的因果量。这一结构因果模型可以表示为描述 $N$ 个单位间因果联系的因果图。在本章的表述中，我们假定网络在基线测量，并在整个随访期间保持静态。同时也指出，这些结果可自然推广到朋友连接集合随时间变化，即时变节点 $F_i(t)$ 本身也具有随机性的情形，例如 van der Laan（2014a）。

如前所述，结构方程模型还对外生误差作了很强的独立性假设；这些假设意味着，给定全部 $N$ 个单位的过去，时点 $t$ 的各单位数据节点在这 $N$ 个单位间相互独立。我们将其称为序贯条件独立假设（sequential conditional independence assumption）。因此，这里假定，时点 $t$ 各单位数据节点之间的任何相依性，都能由全部 $N$ 个单位已观测到的过去完全解释。路线图的下一步，是在明确的附加假设下，从数据分布建立因果量的可识别性；这些假设通常不可检验。可识别性结果使我们能够定义并选定一个包含数据真实概率分布的统计模型，以及一个估计目标，即作用于真实数据分布的目标参数映射；如果所需因果假设成立，这个估计目标就等于该因果量。统计模型必须包含真实数据分布，才能使统计估计目标具有纯统计目标参数的解释；而在为识别因果效应而提出的附加因果条件下，它又可解释为感兴趣的因果量。统计模型，以及将模型中的数据分布映射到参数值的目标参数映射，共同定义了纯统计估计问题。

由于这一统计模型并未假定数据生成实验由独立实验的重复组成，靶向估计量与推断的构造带来了新的统计挑战。靶向最小损失估计（targeted minimum loss-based estimation，TMLE）最初用于独立同分布数据的半参数模型估计（van der Laan 和 Rubin，2006；van der Laan，2008a；van der Laan 和 Rose，2011），随后推广到特定形式的相依处理或删失分配，例如分组序贯自适应设计（van der Laan，2008a；Chambaz 和 van der Laan，2011a,b），以及社区随机试验（van der Laan 等，2013a）。本章介绍靶向最小损失估计在一类新的复杂半参数统计模型中的应用，这些模型允许观测单位之间存在基于网络的个体间干扰。

\section{数据结构}

考虑一项观察 $N$ 个单位样本的研究，其中任意两个单位都可能通过社交网络相连。对单位 $i$，令
\[
 O_i=\bigl(W_i,A_i,\bar Y_i=(Y_i(t):t=1,\ldots,\tau)\bigr)
\]
为按时间排序的纵向观测数据结构，其中 $W_i$ 为基线协变量，$A_i$ 表示行动、处理或暴露干预节点，$Y_i(t)$ 表示单位 $i$ 的时变结局过程，$Y_i(\tau)$ 表示其最终结局。例如，可假定 $A_i$ 记录研究对象 $i$ 的某项特征，如“高水平身体活动”，而 $Y_i(\tau)$ 是该对象在时点 $\tau$ 是否超重的指示变量。令 $F_i$ 为 $W_i$ 的一个分量，并假定 $F_i$ 表示在随访期间任何时点，个体 $i$ 可能从中获得输入的朋友集合。因此，$F_i\subset\{1,\ldots,N\}$。我们假定 $W_i$ 先于 $A_i$ 出现，含义是 $A_i$ 受到包括 $W_i$ 在内的全部过去的影响。

定义 $\mathbf W=(W_i:i=1,\ldots,N)$、$\mathbf A=(A_i:i=1,\ldots,N)$；类似地，对每个 $t=1,\ldots,\tau$，定义 $\mathbf Y(t)=(Y_i(t):i=1,\ldots,N)$。最后，将 $N$ 个单位的观测数据定义为 $\mathbf O=(O_1,\ldots,O_N)\sim P_0^N$。注意，这些观测数据可表示为一个按时间排序的数据结构：
\[
 \mathbf O=\bigl(\mathbf W,\mathbf A,\bar{\mathbf Y}=(\mathbf Y(t):t=1,\ldots,\tau)\bigr).
\]
后一种排序才是唯一具有因果意义的排序。同一时点内各单位的排序由使用者指定，但无关紧要。分别定义 $\operatorname{Pa}(\mathbf A)=\mathbf W$ 和 $\operatorname{Pa}(\mathbf Y(t))=(\mathbf W,\mathbf A,\bar{\mathbf Y}(t-1))$，作为 $\mathbf A$ 及各个 $\mathbf Y(t)$ 的父节点。$A_i$ 的父节点记为 $\operatorname{Pa}(A_i)$，定义为等于 $\operatorname{Pa}(\mathbf A)$；类似地，对 $t=1,\ldots,\tau$ 和 $i=1,\ldots,N$，$Y_i(t)$ 的父节点记为 $\operatorname{Pa}(Y_i(t))$，也定义为等于 $\operatorname{Pa}(\mathbf Y(t))$。

\section{例子}

为使读者了解这一建模方法所适用的现实问题，我们给出几个例子。考虑一项评估暴露前预防（pre-exposure prophylaxis，PrEP）对目标人群，例如 HIV 高风险人群，新发 HIV 感染率之有效性的研究。研究者对一组个体随访数年；对每个个体 $i$，采集其基线特征 $W_i$，其中包含其性伴侣集合 $F_i\subset W_i$。此外，还测量 $i$ 在研究开始时获分配的处理状态 $A_i$，即是否接受 PrEP 的指示变量。假定每隔一定时间定期检测新发 HIV 感染，并以 $Y_i(t)$ 表示单位 $i$ 在时点 $t$ 的 HIV 状态。令 $t=1,\ldots,\tau$，其中第 $\tau$ 个时点表示随访结束。假定我们感兴趣的是，基线 PrEP $A_i$ 对最终时点 $\tau$ 的 HIV 感染比例 $N^{-1}\sum_iY_i(\tau)$ 的效应。

在这一情形中，我们可能希望估计：对 $\mathbf A=(A_1,\ldots,A_N)$ 施加某个随机干预 $\mathbf g^*$ 后的平均结局 $N^{-1}\sum_iE_0Y_{i,\mathbf g^*}(\tau)$。例如，该干预可以将研究中的每个个体都分配为接受 PrEP；这是确定性干预的一个例子。最后，模型中可以假定，给定所有个体的过去，$Y_i(t)$ 和 $A_i$ 的条件分布，除依赖于研究对象 $i$ 自身的过去外，只依赖于其性伴侣各自的过去。换言之，我们假定个体 $i$ 在时点 $t$ 的 HIV 感染风险，取决于其性伴侣的处理状态。

作为随机干预的一个例子，我们可以聚焦于某个特定个体子集，从全部干预节点 $(A_i:i=1,\ldots,N)$ 中选取一个随机子集进行靶向。也就是说，可以定义这样一种随机干预：使某些对象保留其观测暴露，同时将另一些对象的暴露设为预先指定的干预值。例如，若资源有限，只能为少数对象分配处理，研究者可能希望识别出哪些处理分配策略能使感染率下降最多。此外，研究者还可能希望衡量这样一种干预的有效性：作为预防措施，对随机比例的高风险个体给予处理。后一种是真正的随机干预的例子。类似地，也可以将处理节点定义为安全套使用指示变量，使反事实平均结局评估安全套使用对 HIV 流行传播的效应。还可考虑直接干预 $F_i$，例如减少性伴侣数量 $|F_i|$。

再考虑另一项潜在应用：研究单个时点的身体活动对未来超重风险的效应。在时点 $t=1,\ldots,\tau$ 测量个体 $i$ 的超重状态，记为 $Y_i(t)$；其中，最终的 $Y_i(\tau)$ 是个体 $i$ 感兴趣的结局。个体 $i$ 的朋友具有较高身体活动水平，可能产生鼓励作用，进而影响 $i$ 未来的体重，这就是个体间干扰。类似地，个体 $i$ 的身体活动水平也会受到其朋友在基线时体重的影响，这一概念称为自适应处理分配（adaptive treatment allocation）。在这一应用中，个体的处理节点与结局节点，都受到其朋友已观测过去的影响。
